\section{Introduction}
\label{sec:intro}

\subsection{Protein-coding orthology as a tractable problem}
Orthology inference for protein-coding genes is well established and supported by
scalable annotation transfer pipelines
\citep{kirilenkoIntegratingGeneAnnotation2023,sharmaIncreasedAlignmentSensitivity2017,kentEvolutionsCauldronDuplication2003,sharmaCESAR20Substantially2017,herreroEnsemblComparativeGenomics2016,fitchDistinguishingHomologousAnalogous1970,kooninOrthologsParalogsEvolutionary2005}.
Strong sequence conservation, together with constraints imposed by open reading
frames, enables reliable cross-species mapping of gene loci
\citep{sharmaIncreasedAlignmentSensitivity2017,sharmaCESAR20Substantially2017}.
Methods such as TOGA, which combine genome alignment chains with coding sequence
structure, achieve high accuracy at genome scale. However, these approaches rely
on properties specific to protein-coding genes and do not generalize to
non-coding RNAs. In particular, they assume conserved coding structure and
detectable sequence similarity, which are often absent outside protein-coding
contexts.

In this work, we do not attempt to infer evolutionary orthology. Instead, we
identify candidate cross-species correspondences supported by genomic synteny
and embedding similarity, and frame the task as synteny-constrained annotation
transfer.

\subsection{The ncRNA orthology challenge}
Cross-species correspondence is difficult to establish for ncRNAs because the
relevant conserved unit is often unclear. For many ncRNAs, evolutionary
conservation is less readily captured by primary-sequence identity than for
protein-coding genes: related RNAs may retain similar molecular functions,
structural features, or localized conserved elements despite substantial
sequence divergence
\citep{ulitskyLincRNAsGenomicsEvolution2013,
  diederichsFourDimensionsNoncoding2014,
  hezroniPrinciplesLongNoncoding2015,
  rivasEvolutionaryConservationRNA2021,
  ulitskyEvolutionRescueUsing2016,
  statelloGeneRegulationLong2021}.
Conservation may also be restricted to short motifs or local domains rather than
extending across the complete transcript. At the same time, transcript
boundaries and isoform structures are frequently ambiguous, and ncRNA
annotations remain incomplete and heterogeneous across species
\citep{mattickLongNoncodingRNAs2023,perteaCHESSNewHuman2018}.
Together, these properties complicate both annotation transfer and its
evaluation.

These problems are especially acute for lncRNAs. Annotated loci may span tens of
kilobases and contain localized candidate domains separated by long intervening
sequence. Whole-locus comparison can therefore dilute localized correspondence,
and complete loci may exceed the input length supported by current RNA
foundation models. This motivates searching for candidate conserved subregions
rather than requiring end-to-end transcript correspondence.

\subsection{Limitations of existing approaches}

Existing methods address several related tasks, but none directly combines
reference-guided annotation transfer with localization and comparison of
candidate ncRNA subregions within syntenic genomic loci.

Covariance-model approaches implemented in Infernal and distributed through
Rfam can sensitively identify members of established RNA families
\citep{nawrockiInfernal11100fold2013,
  ontiveros-palaciosRfam15RNA2025}. Their application, however, depends on a
predefined family alignment and consensus sequence--structure model. They are
therefore well suited to annotating known RNA families, but less applicable to
reference loci for which no curated family model exists.

Structure-aware alignment methods such as LocARNA and RNAforester compare
supplied RNA sequences using explicit sequence--structure models
\citep{willLocARNA20Versatile2024,hochsmannLocalSimilarityRNA2003}. These methods
address the comparison of already defined RNA intervals, rather than locating a
corresponding element within a genome-scale syntenic search region.

Comparative screening methods such as RNAz and EvoFold identify genomic regions
with evidence of evolutionarily conserved RNA secondary structure in multiple
sequence alignments
\citep{gruberRNAz20Improved2010,pedersenIdentificationClassificationConserved2006}. Their objective is de novo
discovery of candidate structured RNA elements, rather than transfer of a
specific reference annotation to a query genome. They also depend on detectable
comparative sequence and structural signals within the input alignment.

CURIA addresses a different task. Starting from a reference ncRNA annotation and
pairwise genome alignment chains, it first restricts the search to candidate
syntenic loci and then localizes and ranks corresponding subregions using learned
RNA representations. This formulation does not require a pre-existing RNA family
model, a supplied query transcript, or end-to-end conservation of the annotated
reference locus.

\subsection{Synteny-constrained embedding matching}

An underexplored strategy is to use pretrained RNA-language-model
representations as a localized similarity signal within syntenically constrained
genomic search regions. Genome alignment chains provide candidate regions that
preserve genomic context and therefore indicate where a corresponding ncRNA may
be sought
\citep{zoonomiaconsortiumComparativeGenomicsMultitool2020,
  armstrongProgressiveCactusMultiplegenome2020}.
RNA foundation models such as RNA-FM
\citep{chenInterpretableRNAFoundation2022}, RiNALMo
\citep{penicRiNALMoGeneralpurposeRNA2025}, and related approaches such as RNABERT
\citep{akiyamaInformativeRNABase2022} provide learned representations associated
with RNA family identity, secondary-structure-related properties, and downstream
predictive tasks
\citep{vahabAdvancingNoncodingRNA2025}.
Applying these models as unconstrained genome-wide scanners would be
computationally impractical, whereas restricting comparison to chain-defined
candidate loci makes localized RNA-specific matching feasible. This combination
supports comparison of ncRNA loci in settings where primary-sequence conservation
is weak or heterogeneous.

\begin{figure}[t]
  \centering
  % Generated by analysis/make_figures.py (fig1_embeddings) from RiNALMo
  % embeddings cached in analysis/data/fig1_embeddings.npz.
  \includegraphics[width=\linewidth]{fig1_embeddings.pdf}
  \caption{Illustrative embedding behavior supporting local ncRNA comparison.
  \textbf{(A)} Per-nucleotide embeddings
  (matching PCA, first two components) for two human tRNA-Asn-GTT gene copies
  (\textit{tRNA-Asn-GTT-chr1-139} and \textit{-140}, $74$\,nt each) and one
  pre-miRNA (\textit{MIR103A1}, ENST00000362168, $110$\,nt), all from the human
  genome assembly hg38; each point is a nucleotide. The tRNA copies cluster
  together while the miRNA forms a distinct cluster. \textbf{(B)} Out-of-fold
  probability from a 5-fold cross-validated logistic classifier trained on the
  embeddings to distinguish annotated hg38 ncRNAs (purple; tRNA, snoRNA, miRNA)
  from background sequences (gray; intergenic and dinucleotide-shuffled): ncRNAs
  score near $1$ and background near $0$ (AUC ${\approx}0.96$), showing the
  embeddings separate the selected RNA-derived sequences from genomic background.
  Panel (B), titled ``Illustrative embedding separability'' within the figure, is
  an illustrative separability experiment on a small balanced set and is not the
  production island detector used in the pipeline (Section~\ref{sec:refisland}),
  which is trained on a different biotype set and negatives and operates at a
  different threshold.
  \textbf{(C)} Maximum Mean Discrepancy (MMD) between a short-ncRNA reference and
  windows tiled across its extended locus, versus offset from the annotated
  position (median and interquartile range over $n=100$ short ncRNAs, each locus
  min--max normalized). MMD is minimised at the true position (offset $0$) and
  rises with positional shift. \textbf{(D)} Best-hit MMD versus a distant control
  window ($250$\,nt away) for the same $n=100$ loci, with marginal histograms.
  Best-hit MMD concentrates near $0$ (31 loci essentially perfect), whereas the
  distant control almost never does; 91\% of loci lie above the diagonal
  (matched window more similar than the control).}
  \label{fig:embeddings}
\end{figure}

\subsection{Method overview}

We formulate cross-species ncRNA correspondence as the matching of
embedding-supported regions within loci defined by genome alignment chains.
Candidate loci are first identified using a chain-based projection adapted from
principles used in TOGA. The projection uses alignment coverage across exonic,
intronic, and flanking sequence to provide locus-level syntenic context that is
not determined solely by conservation within the ncRNA itself
\citep{kirilenkoIntegratingGeneAnnotation2023,
  sharmaIncreasedAlignmentSensitivity2017}.

Within each candidate locus, \curia{} compares RNA-language-model
representations using one of two geometry-dependent branches. Compact,
single-exon ncRNAs are compared as complete regions using Maximum Mean
Discrepancy (MMD)
\citep{grettonKernelTwoSampleTest2012}, which measures the difference between
the distributions of their per-position embedding vectors. Longer or
multi-exon loci are processed by a localized-subregion workflow, in which
detector-positive subregions are matched by Smith--Waterman local alignment
\citep{smithIdentificationCommonMolecular1981} over a per-token cosine-similarity
matrix. Both branches provide continuous scores for ranking candidate
correspondences without requiring nucleotide sequence alignment.

\subsection{Contributions}
In this work, we:
\begin{enumerate}
  \item Formulate cross-species ncRNA annotation as synteny-constrained matching
  in RNA embedding space.
  \item Introduce \curia{}, a pipeline that combines chain-based locus projection
  with localized embedding-based comparison.
  \item Develop complementary strategies for compact and extended ncRNA loci,
  including boundary refinement for short ncRNAs and localized subregion
  matching for longer or multi-exon loci.
  \item Demonstrate the approach at genome scale across 19 mammalian genomes and
  identify recurrent candidate correspondences that are not fully explained
  by conventional sequence-similarity measures.
\end{enumerate}
